%! The Quadratic Formula & the Discriminant — Unit 3, Lesson 3.5 — Warm-Up
\documentclass[10pt]{article}
\usepackage{algebra2-article}
\usepackage{algebra2-boxes}
\newcommand{\TallMath}[1]{$\displaystyle #1\rule[-1.4em]{0pt}{3.2em}$}

\begin{document}

\pageheader{Unit 3, Lesson 3.5}{Warm-Up}

\vspace{0.10in}

\begin{notesbox}{Getting Ready: Skills We'll Use Today}
\small These three skills are the machinery of the quadratic formula. Work carefully --- signs matter.

\vspace{4pt}
\textbf{1. Evaluate $b^2-4ac$.} Substitute the given values; watch every sign.
\[
  a=1,\ b=-2,\ c=-3:\ \ b^2-4ac=\blank{2.2cm} \qquad
  a=1,\ b=6,\ c=9:\ \ b^2-4ac=\blank{2.2cm}
\]
\[
  a=2,\ b=3,\ c=5:\ \ b^2-4ac=\blank{2.4cm}
\]

\vspace{4pt}
\textbf{2. Simplify each radical.} Pull out perfect squares; use $i$ for the negatives (Lesson 3.4).
\[
  \sqrt{16}=\blank{1.6cm} \qquad \sqrt{20}=\blank{1.8cm} \qquad
  \sqrt{-16}=\blank{1.8cm} \qquad \sqrt{-8}=\blank{1.8cm}
\]

\vspace{4pt}
\textbf{3. Write in standard form, then identify $a$, $b$, $c$.} Move every term to one side, $=0$.
\begin{itemize}\setlength{\itemsep}{3pt}
  \item $2x^2=5x-1 \ \Rightarrow\ $ \blank{3.0cm} \hfill $a=\blank{1.0cm}\ \ b=\blank{1.0cm}\ \ c=\blank{1.0cm}$
  \item $x^2+9=4x \ \Rightarrow\ $ \blank{3.0cm} \hfill $a=\blank{1.0cm}\ \ b=\blank{1.0cm}\ \ c=\blank{1.0cm}$
\end{itemize}
\emph{Today: one formula, $x=\dfrac{-b\pm\sqrt{b^2-4ac}}{2a}$, solves them all --- and that first quantity you evaluated has a name.}
\end{notesbox}

\end{document}
